\documentclass[11pt]{scrartcl}
\usepackage[utf8]{inputenc}
\usepackage[sexy]{evanchen} 
\usepackage{graphicx}
\graphicspath{{8 Outsource K Math/Gambar/}}

\title{Integers Bounding}
\author{Azzam and Akeyla}

\begin{document}

\maketitle

\renewcommand*\contentsname{Daftar Isi}
\tableofcontents

\newpage

\section{Bounding}
Membatasi nilai sehingga hanya tersisa berhingga kemungkinan nilai. 

\begin{example}
    Ada berapa banyak bilangan asli $x$ dan $y$ yang memenuhi persamaan $x^{10} + y^2 = 2048$? 
\end{example}

\begin{proof}[Solusi]
    Dikarenakan $y^2$ pasti positif, maka $x^{10}$ harus bernilai maksimal $2047$. Namun, jika $x \geq 3$, maka $x^{10} \geq 3^{10} = 59049$ yang jauh lebih besar dari $2048$. Sehingga, $x \leq 2$. Dapat dicek bahwa yang memenuhi hanyalah $x = 2, y = 32$
\end{proof}

Terkadang, berasumsi tanpa mengurangi keumuman diperlukan.

\begin{example}
    Tentukan banyak tripel bilangan asli $(x, y, z)$ yang memenuhi persamaan \[\dfrac{1}{x} + \dfrac{1}{y} + \dfrac{1}{z} = \dfrac{3}{2}\]
\end{example}

\begin{proof}[Solusi]
    Asumsikan tanpa mengurangi keumuman bahwa $x \geq y \geq z$. Artinya, $\dfrac{1}{x} \leq \dfrac{1}{y} \leq \dfrac{1}{z}$ dan $\dfrac{3}{2} = \dfrac{1}{x} + \dfrac{1}{y} + \dfrac{1}{z} \leq \dfrac{3}{z}$. Dapat disimpulkan bahwa $z \leq 2$.

    Selanjutnya, bagi kasus berdasarkan nilai $z$.

    \textbf{Kasus 1: $z = 1$} Maka, $\dfrac{1}{x} + \dfrac{1}{y} = \dfrac{1}{2}$. Dengan mengubah bentuk persamaan tersebut, diperoleh $xy = 2x + 2y$. Bentuk tersebut ekuivalen dengan $(x - 2)(y - 2) = 4$ dan memiliki solusi $(3, 6)$ dan $(4, 4)$. Jadi, diperoleh solusi $(1, 3, 6)$ dan $(1, 4, 4)$ beserta dengan permutasinya.

    \textbf{Kasus 2: $z = 2$} (Kurang lebih sama, lanjutan solusi diserahkan ke pembaca)
\end{proof}

Fungsi-fungsi tertentu tumbuh lebih cepat daripada fungsi-fungsi lainnya sehingga solusi yang diberikan harus berasal dari bilangan yang "cukup kecil".

\begin{example}
    Tentukan semua bilangan asli yang memenuhi $2^x = x^2$
\end{example}

\begin{proof}[Solusi]
    Perhatikan bahwa untuk $x \geq 5$. Nilai $2^x$ akan lebih besar dari $x^2$. Sehingga, haruslah $x \leq 4$. Pengecekan lebih lanjut akan menunjukkan bahwa solusinya hanyalah $x = 2$ dan $4$.
\end{proof}

% Terakhir, kita juga dapat menggunakan berbagai teknik ketaksamaan aljabar. Dua ketaksamaan yang umum digunakan untuk menyelesaikan persamaan adalah QM-AM-GM-HM dan Cauchy-Schwarz.

% \begin{theorem}[Ketaksamaan QM-AM-GM-HM]
% Untuk $n$ bilangan asli dan $a_1,a_2,a_3,...,a_n$ bilangan real positif, notasikan \\
% QM = $\sqrt{\dfrac{a_1^2 + a_2^2 +... + a_n^2}{n}}$ AM = $\dfrac{a_1 + a_2 + ... + a_n}{n}$ GM = $\sqrt[n]{a_1a_2...a_n}$ \newline HM = $\dfrac{n}{\dfrac{1}{a_1} + \dfrac{1}{a_2} + ... \dfrac{1}{a_n}}$\\
% Maka, $$QM \geq AM \geq GM \geq HM$$
% dengan kesamaan terjadi ketika $a_1 = a_2 = ... = a_n$
% \end{theorem}

% \begin{theorem}
%     [Ketaksamaan Cauchy-Schwarz] Untuk setiap bilangan real $a_1,a_2,...,a_n$ dan $b_1,b_2,...,b_n$, berlaku $$(a_1^2 + a_2^2 + ... + a_n^2)(b_1^2 + b_2^2 +... + b_n^2) \geq (a_1b_1 + a_2b_2 + ... + a_nb_n)^2$$ dengan kesamaan terjadi jika salah satu dari tiga hal berikut terpenuhi : 
%             \begin{enumerate}
%                 \item $a_1 = a_2 = ... = a_n = 0$
%                 \item $b_1 = b_2 = ... = b_n = 0$
%                 \item Terdapat suatu bilangan real $k$ sedemikian hingga $\frac{a_1}{b_1} = \frac{a_2}{b_2} = ... = \frac{a_n}{b_n} = k$
%             \end{enumerate}
% \end{theorem}

% \begin{example}[Beberapa contoh akibat]
%     Untuk bilangan asli $x, y, z$ berlaku
%     \begin{enumerate}
%         \item $x^2 + y^2 \geq 2xy$
%         \item $x^3 + y^3 + z^3 \geq 3xyz$
%         \item $\dfrac{x}{y} + \dfrac{y}{z} + \dfrac{z}{x} \geq 3$
%         \item $(x + y)\left(\dfrac{1}{x} + \dfrac{1}{y}\right) \geq 4$
%     \end{enumerate}
% \end{example}

\section{Soal-soal}
    \begin{problem}
        Diberikan bilangan cacah $a, b, c$ yang memenuhi persamaan $120 = 7a + 19b + 31c$. Hitunglah nilai dari $a + b + c$.
    \end{problem}

    \begin{problem}
        Carilah semua bilangan bulat $a, b$ yang memenuhi persamaan \[6a^2 + 5b^2 = 74\]
    \end{problem}
    
    \begin{problem}
        Tentukan semua bilangan asli $x$ yang memenuhi persamaan \[\dfrac{1}{x} + \dfrac{1}{x + 1} + \dfrac{1}{x + 2} = \dfrac{13}{12}\]  
    \end{problem}

    \begin{problem}
        Carilah semua solusi bilangan asli $p$ yang memenuhi \[(p-2)^3 = 18p-1\]
    \end{problem}

    \begin{problem}
        Banyak bilangan asli $a, b, c$ yang memenuhi persamaan \[\dfrac{1}{a} + \dfrac{1}{b} + \dfrac{1}{c} = \dfrac{3}{5}\] adalah ...
    \end{problem}

    % \begin{problem}
    %     Tentukan solusi bilangan asli yang memenuhi sistem persamaan \[a + b + c + d = 12\] dan \[abcd = 27 + ab + ac + ad + bc + bd + cd\]
    % \end{problem}

    \begin{problem}
        Banyak tripel bilangan asli $a, b, c$ yang memenuhi persamaan \[abc = a + b + c\] adalah ...
    \end{problem}

    % \begin{problem}
    %     Tentukan semua tupel bilangan asli $(x, y, z, u, v)$ yang memenuhi persamaan \[xyzuv = x + y + z + u + v\]
    % \end{problem}

    % \begin{problem}
    %     Tentukan semua bilangan prima $a, b, c$ yang memenuhi \[ab + bc + ca > abc \]
    % \end{problem}

    \begin{problem}
        Carilah semua pasangan bilangan bulat $x, y$ yang memenuhi persamaan \[x^3 + y^3 = (x + y)^2\]
    \end{problem}

    % \begin{problem}
    %     Carilah semua solusi bilangan asli $m,n$ yang memenuhi persamaan \[(m+1)! + (n+1)! = m^2n^2\]
    % \end{problem}

    % \begin{problem}
    %     Tentukan nilai bilangan asli $n, k_1, k_2, ..., k_n$ yang memenuhi sistem persamaan \[k_1 + k_2 + ... + k_n = 5n - 4\] dan \[\dfrac{1}{k_1} + \dfrac{1}{k_2} + ... + \dfrac{1}{k_n} = 1\]
    % \end{problem}

    \begin{problem}
        Tentukan semua bilangan asli $a, b, c$ yang memenuhi persamaan \[\left(1 + \dfrac{1}{a}\right)\left(1 + \dfrac{1}{b}\right)\left(1 + \dfrac{1}{c}\right) = 2\]
    \end{problem}

    % \begin{problem}
    %     Tentukan semua bilangan prima $p$ dan $q$ sehingga $p$ membagi $3q - 1$ dan $q$ membagi $3p - 1$
    % \end{problem}
\end{document}